% =====================================================================
% Results -- query-conditioned routing: keyword router vs LLM router,
% and one query traced through the whole graph.
%
% Paste (or \input) this subsection into the paper's results section.
% Preamble needs:  \usepackage{booktabs} \usepackage{array}
%
% All three tables are WRITTEN BY THE NOTEBOOK (test 16D) -- never edit them:
%   tables/tab_routing_summary.tex   keyword vs LLM totals      (one column)
%   tables/tab_routing_plain.tex     the twelve plain prompts   (table*, two columns)
%   tables/tab_prompt_flow.tex       two queries, stage by stage (table*, two columns)
%
% Numbers quoted in the prose come from notebook tests 15C, 16B and 16C and
% from experiments/routing/llm_routing_runs.json. If anything is re-run,
% re-check them: 29 prompts, 12/12, 0/12, 9/12, 14/29, 2/29, 1.79, 1.41,
% 29/29, 3.0 s, 2.8-3.6 s, 87 runs, 100 vs 80 deliveries.
%
% If space is short, tab_prompt_flow can move to an appendix; the summary
% table and the plain-prompt table carry the result.
% =====================================================================

\subsection{Query-Conditioned Routing}
\label{sec:routing}

% Tables are defined here, in the order the text first cites them, so they
% are numbered in that order. The wide ones (table*) can only appear at the
% top of a page AFTER the one they are defined on, so they must come early.

\begin{table}[t]
  \centering
  \footnotesize
  \setlength{\tabcolsep}{5pt}
  \caption{Keyword router versus LLM router on the 29 prompts. A target counts
    as activated only when the router chose it, not when it was reached through
    the default pair.}
  \label{tab:routing-summary}
  % Generated by notebooks/cricketmind_langgraph_demo.ipynb (test 16D).
% Do not edit by hand -- re-run the notebook instead.
\begin{tabular}{@{}lrr@{}}
\toprule
 & Keyword & LLM \\
\midrule
Keyword prompts: target activated & 12/12 & 12/12 \\
Plain prompts: target activated & 0/12 & 9/12 \\
Fell back to the default agents & 14/29 & 2/29 \\
Agents activated per prompt (mean) & 1.79 & 1.41 \\
Identical routing in all 3 runs & -- & 29/29 \\
Routing time per prompt (mean) & -- & 3.0\,s \\
\bottomrule
\end{tabular}

\end{table}

\begin{table*}[t]
  \centering
  \footnotesize
  \caption{The twelve plain prompts, which avoid every trigger word of their
    target agent, and the agents the LLM router activated. \textbf{Bold}: the
    target was activated by the router's own choice. The keyword router
    activated the default pair (Coaches, Players) for all twelve.}
  \label{tab:routing-plain}
  % Generated by notebooks/cricketmind_langgraph_demo.ipynb (test 16D).
% Do not edit by hand -- re-run the notebook instead.
\begin{tabular}{@{}l>{\raggedright\arraybackslash}p{0.52\textwidth}>{\raggedright\arraybackslash}p{0.30\textwidth}@{}}
\toprule
Written for & Plain question (no trigger words) & LLM router activated \\
\midrule
Franchises & Should our team keep Kohli next season, or let him go to save money? & \textbf{Franchises}, Players, Player Agents \\
Coaches & How do we get Kohli out early when he faces spin? & Players \\
Players & Has Kohli's batting got better or worse this season, and what should he change? & \textbf{Players} \\
Scouts & Find young overseas batters who could become the next Kohli. & \textbf{Scouts}, Players \\
Fans & I want to watch Kohli bat live; which home game should I go to? & Players, Franchises \\
Stadiums & How should the ground staff get the wicket ready if we want Kohli to struggle? & \textbf{Stadiums}, Medical Teams \\
Sponsors & Is paying to put our company name on Kohli's bat good value for money? & \textbf{Sponsors} \\
Broadcasters & What on-screen visuals should the TV team show when Kohli walks out to bat? & \textbf{Broadcasters} \\
Governing Bodies & Did Kohli break any of the game's official regulations in this innings? & \textbf{Governing Bodies} \\
Academies & How can we teach teenage batters to hit the ball through the covers like Kohli? & Players \\
Medical Teams & Is Kohli too tired or sore to play the final safely? & \textbf{Medical Teams}, Players \\
Player Agents & How much money should Kohli ask for in his next deal? & \textbf{Player Agents} \\
\bottomrule
\end{tabular}

\end{table*}

\begin{table*}[t]
  \centering
  \footnotesize
  \caption{Two queries traced through the graph. A: \emph{``Compare Kohli's
    cover drive vs his pull shot against fast bowlers in the powerplay''}.
    B: \emph{``Is Kohli's workload a fatigue risk before the final?''}. Each
    cell is what the stage wrote to the shared state; agent cells list every
    figure the agent cites. Stages whose output is identical for both queries
    (perception, Shot DNA, context retrieval, output) are omitted.}
  \label{tab:prompt-flow}
  % Generated by notebooks/cricketmind_langgraph_demo.ipynb (test 16D).
% Do not edit by hand -- re-run the notebook instead.
\begin{tabular}{@{}>{\raggedright\arraybackslash}p{0.18\textwidth}>{\raggedright\arraybackslash}p{0.38\textwidth}>{\raggedright\arraybackslash}p{0.38\textwidth}@{}}
\toprule
Stage & Question A & Question B \\
\midrule
N3 reads the question & bowler\_type=pace, phase=powerplay, player=Kohli, shots=Cover Drive+Pull & player=Kohli \\
N1-N5 classify the clip & Cover Drive (p=0.92) & Cover Drive (p=0.92) \\
N6 builds the Shot DNA & Cover Drive + posture (weight transfer 0.88) + filters & Cover Drive + posture (weight transfer 0.88) + filters \\
N7a finds deliveries & 100 matching deliveries & 80 matching deliveries \\
N7b finds context & 4 context rows & 4 context rows \\
N8 computes statistics & Cover Drive SR 142.0 (50 balls), Pull SR 112.0 (50 balls) & Cover Drive SR 130.0 (80 balls) \\
N9 simulates & 6 consecutive deliveries targeting the pull: 39.4\% out (95\% CI 38.7\%-40.1\%) & 6 consecutive deliveries targeting the cover drive: 13.9\% out (95\% CI 13.4\%-14.4\%) \\
N10 picks the experts & wakes Players, Coaches & wakes Medical Teams \\
N11 agents answer & Coaches: best\_sr=142.0, worst\_sr=112.0, p\_pct=39.4; Players: best\_sr=142.0, worst\_sr=112.0, weight\_transfer=0.88 & Medical Teams: sample\_size=80 \\
Faith Check & all cited figures match the engines & all cited figures match the engines \\
N13 output & status = ok & status = ok \\
\bottomrule
\end{tabular}

\end{table*}

We compare the keyword router and the language-model router of
Section~\ref{sec:method} on the same queries. The language model is Qwen2.5-7B,
served locally through Ollama with temperature~0 and a fixed seed; for a fair
comparison it receives exactly what the keyword router knows, each agent's
stakeholder name and keyword list. The instructions were fixed before the
comparison was run, and the same rules follow both proposals, so the bounds of
Section~\ref{sec:exec-profile} hold for either router.

\paragraph{Prompt set}
For each of the twelve stakeholders we wrote two questions asking for the same
thing: a \emph{keyword} version that uses the agent's trigger words (``Is Kohli's
workload a fatigue risk before the final?'') and a \emph{plain} version that
avoids all of them (``Is Kohli too tired or sore to play the final safely?'').
Five further prompts cover a question with no signal, one naming six
stakeholders at once, the running example, and two text-only questions---29 in
all. The notebook verifies that every keyword prompt contains its target's
trigger words and that no plain prompt contains any. Each prompt was routed
once by the keyword router, which is deterministic, and three times by the LLM
router.

\paragraph{Results}
Table~\ref{tab:routing-summary} summarises the comparison. Both routers
activate the intended agent for all twelve keyword prompts. On the plain
prompts the keyword router activates none of the twelve: finding no trigger
word, it falls back to the default pair on every one, and on 14 of the 29
prompts overall. The LLM router activates the intended agent for nine of the
twelve plain prompts (Table~\ref{tab:routing-plain}) and falls back only twice.
Both fall-backs are text-only prompts for which the model proposed only agents
that cite a posture measurement; with no clip there is no posture, so the
evidence rule removed them. The LLM router is also more selective, activating
1.41 agents per prompt against 1.79. On the prompt naming six stakeholders, the
keyword router's four-agent cap, with ties broken by registry order, discarded
Sponsors, Medical Teams and Broadcasters; the LLM router kept all three.

The LLM router's answers were identical in all three runs for all 29 prompts,
and a routing call took 3.0\,s on average (2.8--3.6\,s) on a single consumer
GPU. Across all 87 LLM-routed runs, no activated agent lacked its evidence and
no run exceeded the cap; this holds by construction, since the rules follow the
proposal, and the runs confirm it.

\paragraph{Where the LLM router fails}
The three plain prompts it missed are informative. Asked how to dismiss Kohli
cheaply against spin, it chose the batter's perspective (Players) instead of
the coaching staff's; asked which home game a supporter should attend, and how
to teach teenage batters a stroke, it again chose Players rather than Fans or
Academies. The router also sees only the query text. For the running example
the keyword router additionally activated Coaches because of a 30-point
strike-rate spread in the analytics, a signal the language model never
receives.

\paragraph{One query through the graph}
Table~\ref{tab:prompt-flow} traces two queries through every stage, recorded
from the framework's update stream rather than written by hand. The same
pipeline extracts different filters, retrieves different evidence (100 against
80 deliveries), computes different statistics, and activates different agents,
and every figure an agent cites is checked against the engines before release.
The trace also exposes a cost: for the medical query, the strike rates and the
dismissal simulation are computed although the activated agent cites only the
sample size, because every analytics stage runs before routing.

% The limitations of this comparison (author-written prompts, twelve plain
% prompts, single-machine stability) are stated once, in limitations.tex.
